\chapter*{Preface}
\addcontentsline{toc}{chapter}{Preface}
\markboth{Preface}{Preface}

In 1967 Robert Langlands wrote a letter to Andr\'e Weil proposing a web of conjectures that would connect automorphic forms, the representation theory of reductive groups, and the arithmetic of Galois representations. Six decades later the Langlands programme is a central organising project of number theory. Many of its cases are theorems, among them the modularity of elliptic curves over the rationals, the local Langlands correspondence for $\GL_n$, the fundamental lemma, and symmetric power functoriality for holomorphic modular forms. The general conjectures remain open.

This book asks what probability can contribute. Our answer has two sides.

The first is a limitation. No probabilistic argument known to us proves reciprocity or functoriality, and in Chapter~\ref{ch:01} we prove a proposition explaining why no argument that sees only the limiting distribution of Frobenius data could. Equidistribution is a consequence of functoriality, not a route to it. We do not claim to prove the Langlands conjectures or any Riemann hypothesis.

The second is what probability does supply: consequences of the conjectures, calibration against data, and probabilistic reformulations. The Langlands conjectures predict that the Frobenius conjugacy classes attached to an automorphic representation behave like Haar random elements of a compact group. That prediction is the probabilistic shadow of the programme.

Part~II is about equidistribution: Weyl's criterion and the Sato Tate law, Frobenius statistics at scale, random matrix models of the zeros of $L$ functions, Katz's vertical law for Kloosterman sums and the random Fourier series of Kowalski and Sawin, and the Cohen Lenstra and Poonen Rains heuristics, in which class groups and Selmer groups behave like random groups.

Part~III treats Hecke operators as random walks on Bruhat Tits trees, Ihara zeta functions, Ramanujan graphs, and hitting times on supersingular isogeny graphs.

Part~IV starts from the infinitely divisible laws attached to Euler products and follows Selberg's central limit theorem to the log correlated picture of the critical line, with its branching random walk, multiplicative chaos and the conjecture of Fyodorov, Hiary and Keating. At its centre is the Thorin framework, which extends earlier work of the first author \cite{polson2018vdp}. For every cuspidal automorphic representation of $\GL_n$ over $\Q$, the generalized Riemann hypothesis is equivalent to the statement that the reciprocal of an even completion of its $L$ function, read in the variable $s=z^2$, is the Laplace transform of a generalized gamma convolution; equivalently, that a L\'evy function built from the zeros is completely monotone, that a logarithmic derivative is a Stieltjes function, or that a sequence computable from central derivatives alone is a Stieltjes moment sequence. These are reformulations, not proofs. The same structure appears locally, where the Ramanujan conjecture at a prime is exactly a local Thorin property; for zeta functions of curves over finite fields, where Weil's theorem makes it unconditional; on supersingular isogeny graphs, where the weight two bound of Eichler, Shimura and Weil does; and for Hecke random walks on Bruhat Tits trees, whose passage times are generalized gamma convolutions with Thorin measure filling the Kesten band. The framework also connects to van Dantzig pairs and Wald couples, to the sine squared transport that carries the primes, through the explicit formula, to the Thorin measure, and to the work of Konstantopoulos, Patie and Sarkar on the van Dantzig problem, where we give a necessary condition for their class and test it on fourteen $L$ functions. Chapter~\ref{ch:squarefree} reads the error term of the squarefree counting function through the P\'olya law of $\xi$ and its Thorin atoms; Chapter~\ref{ch:beta} carries the framework to Dirichlet's $\beta$ function, whose P\'olya kernel is built from the cube of Dedekind's eta function, and proves, up to the heights to which the Riemann hypotheses for $\zeta$ and $\beta$ are verified, an inequality on the lines $\frac12<\operatorname{Re}s\le3.15$ that Lagarias proved for $\xi$ on every line to the right of the centre under the Riemann hypothesis; Chapter~\ref{ch:scholze} computes signed Thorin densities by Stieltjes inversion and counts zeros numerically by the argument principle for $\zeta$, $L(s,\chi_{12})$ and $L(s,\Delta)$. None of the three proves a Riemann hypothesis. Chapter~\ref{ch:recip} turns these devices on Langlands reciprocity: it derives quadratic reciprocity from the normal law, treats Frobenius traces as random variables and as the labels of a Gamma process, and recasts Hecke's converse theorem at level one as a symmetry of a P\'olya kernel, and the holomorphy and the generalized Riemann hypothesis of a self dual Artin $L$ function as the complete monotonicity of a heat trace computed from Frobenius traces, the conductor and the gamma factor. Again these are reformulations, and no new case of reciprocity is proved. Chapter~\ref{ch:hit} reads reciprocity laws as identities between hitting times: on supersingular isogeny graphs in characteristic $N$, Kemeny constants and, for $N\equiv1\bmod12$, numbers of spanning trees are computed from the reduction of $J_0(N)$ modulo the degree of the isogenies, Chebotarev densities are hitting rates for the largest prime factors of a random integer, the Thorin converse theorem concerns explosion times of pure birth processes, and the functional equation of the $L$ function of an odd primitive Dirichlet character is an identity between two expansions of a signed exit law of Brownian motion. Here too no new case of reciprocity is proved. Chapter~\ref{ch:flt} follows the proof of Fermat's last theorem through these devices: weight one forms, through which the theorem of Langlands and Tunnell gives the residual modularity that Wiles's argument starts from, are the newforms whose label laws are supported on finitely many points; the modularity of an elliptic curve is, by Weil's converse theorem, a family of reflection symmetries of twisted P\'olya kernels built from point counts and the conductor; Mazur's torsion theorem, through an argument of Serre, makes positive, for a Frey curve and $p\ge5$, the Gamma shape of the primes $\ell$ at which $p$ does not divide $|E(\mathbb F_\ell)|$, the local density of Koblitz's conjecture; and the genus zero of $X_0(2)$ empties a Thorin measure built from Hecke eigenvalues at level $2$. These restate the proof, reprove no step of it and prove no new case of modularity; its content is the theorems of Langlands and Tunnell, Mazur, Ribet, Wiles, and Taylor and Wiles, on the curve of Hellegouarch and Frey and the reduction of Frey and Serre. Chapter~\ref{ch:lap} collects the inverse Laplace transforms behind these devices: Bohr means read Frobenius traces off the transported density of $\eta(z)\eta(23z)$; Bromwich inversion of the Artin $L$ function of $x^3-x-1$ on the critical line returns that weight one form, given the holomorphy that Hecke's identification of the two $L$ functions supplies; moving the basepoint adds Cauchy noise to the zeros; the Thorin variable of an Euler product with nonnegative coefficients has a Pareto tail of index $\frac32$, as Roynette and Yor found; atomic Thorin laws invert to B splines; and the Sato Tate law minimises the roughness penalty of Good and Gaskins among densities that vanish on the walls of the alcove. Most of these results are known, and none of them proves a case of reciprocity or a Riemann hypothesis. Chapter~\ref{ch:pd} relates the Gamma process of Chapter~\ref{ch:recip} to classical laws: the normalised sizes of the prime factors of a random integer follow the Poisson Dirichlet law, as the normalised cycle lengths of a random permutation do; normalised, the Gamma process with Frobenius labels is Ferguson's Dirichlet process, with the Haar measure of Chebotarev's theorem as base measure; and the random integers weighted by the coefficients of the powers of $\zeta$ give the whole Ewens family. For uniform random integers below $x$, the class of a size biased prime factor approaches Haar measure at rate $1/\log x$, the class of the largest one faster still, but the classes of the distinct prime factors only at rate $1/\log\log x$, with Mertens constants per Chebotarev class computed from $L(1,\chi)$ and $\frac{L'}L(1,\chi)$. These are theorems about random integers, most of them known, and none proves a case of reciprocity. Chapter~\ref{ch:gl1} treats degree one, where reciprocity is a theorem: Dirichlet's theorem on primes in progressions and Hecke's theorem on the angles of Gaussian primes become Dirichlet processes, Landau's theorem on the integers with no prime factor in a given class becomes an atom at zero, and for an elliptic curve with complex multiplication, whose $L$ function is that of a Hecke character by Deuring's theorem, the Gamma shape test of Chapter~\ref{ch:recip} detects the complex multiplication. Chapter~\ref{ch:maass} turns to the Maass forms, the other half of the automorphic spectrum of $\SL_2(\Z)$: it reads the Selberg trace formula as a heat trace, whose terms are return probabilities of Brownian motion on a hyperbolic surface, restates Selberg's eigenvalue conjecture as the complete monotonicity of that heat trace, treats the closed geodesics as primes to which the sine squared transport and the Poisson Dirichlet limits apply, and reads the Eichler Selberg trace formula as the moment identity behind Serre's vertical Sato Tate law. These too are reformulations, and no new case of Selberg's conjecture is proved. Chapter~\ref{ch:percolation} takes up the other body of Langlands's work, on crossing probabilities in critical percolation: Cardy's formula is a Beta law at the modular lambda function, whose log odds is an extended generalized gamma convolution, and the aspect ratio at which a critical rectangle stops being crossed has, up to a constant, the completed $L$ function of the elliptic curve $y^2=x^3+1$ as its Mellin transform, so that the generalized Riemann hypothesis for that curve becomes a positivity statement about cumulants of a crossing probability. Here too no Riemann hypothesis is proved. Chapter~\ref{ch:cal} reads the theorems of Calegari and his coauthors on even Galois representations, potential modularity and torsion in the dictionary of the book: the comparison of the two gamma factors of the Galois heat trace holds in every weight, the Dirichlet process with the Sato Tate base measure is a theorem over every CM field, and the Thorin criterion becomes a well posed statement for abelian surfaces. These are consequences of known theorems. Part~IV ends with an extreme value representation of the $L$ function of a multiplicative eta product as the expectation of a regularised incomplete gamma function at a maximum of independent exponential variables.

Part~V turns to inference: right Haar measure, transformation models and Barndorff Nielsen's $p^*$ formula, reference priors for orbits after Chang and Eaves, equidistribution of Heegner points, and Bayesian identification of Sato Tate groups. Chapter~\ref{ch:13} states twelve open problems.

Results new to this book include the limit $\ell/(\ell-1)$ for Kemeny constants of isogeny graphs, the Thorin criteria for the generalized Riemann hypothesis of automorphic $L$ functions and for the Ramanujan conjecture, the Thorin structure of Hecke passage times, the representation of the squarefree error under Gamma($\nu$) smoothing as an average of the Mertens function, with positive weights exactly when $\nu\le2$, the signs of the P\'olya kernels of imaginary quadratic fields, the minimum on the real axis of the logarithmic derivative of the completed Dirichlet $\beta$ function on the lines $\frac12<\operatorname{Re}s\le3.15$, up to the height to which its Riemann hypothesis is verified, the Gamma process limit of the Frobenius labelled factorisation of a random integer, the expansion in $1/\log x$ of the probability that a size biased prime factor of a uniform random integer below $x$ has Frobenius in a given conjugacy class, the extreme value representation, and the consistency theorem for Bayesian identification.

Computation runs through the book as the experimental side of the subject. We computed Frobenius traces for all primes below $10^7$, Hecke spectra on supersingular isogeny graphs at 207 prime levels, zeros of $L(s,\Delta)$, and posterior distributions over Sato Tate groups of genus two curves. Every chapter before Chapter~\ref{ch:14} ends with open computational problems, stated with their scale and expected deliverables but deliberately left unsolved. Chapter~\ref{ch:14} collects them into a programme, the agenda for the second author's computational work and, we hope, for some of our readers. Appendix~\ref{app:B} prints the core routines.

The book is written for probabilists and statisticians who want to see where their tools meet the Langlands programme, and for number theorists curious about what probabilistic structure lies inside their objects. We assume graduate probability and some familiarity with modular forms; Chapter~\ref{ch:02} and Appendix~\ref{app:A} fill the gaps on each side.

\bigskip
\noindent Nick Polson, Chicago\\
Vadim Sokolov, Fairfax\\
October 2026
